% Answers to Chapter 13, from lectures/L13/appendix/c_solutions.md; the self-test from
% lectures/L13/README.md.

\facitavsnitt{c13}{Kapitel 13}{Derivata, del II}
\begin{facit}

\svar{1.1} \sv{a} $f'(x) = -1{,}5x^2 + 3x + 4{,}5$ \sv{b} $x = -1$ och $x = 3$
\sv{c} Maximum i $x = 3$ ($f''(3) = -6 < 0$), minimum i $x = -1$ ($f''(-1) = 6 > 0$).
\sv{d} Maximivärdet är $f(3) = 10{,}5$ och minimivärdet $f(-1) = -5{,}5$.
\sv{e} Kurvan har en dal i $(-1;\, -5{,}5)$ och en topp i $(3;\, 10{,}5)$ och skär $y$-axeln i
$-3$.

\svar{1.2} \sv{a} $u'(t) = -0{,}6t^2 + 2{,}4t - 1{,}4625$
\sv{b} $t_1 = 0{,}75\,\text{s}$ och $t_2 = 3{,}25\,\text{s}$
\sv{c} $u''(0{,}75) = 1{,}5 > 0$: minimum; $u''(3{,}25) = -1{,}5 < 0$: maximum.
\sv{d} Minimivärdet är $u(0{,}75) \approx 2{,}49\,\text{V}$ och maximivärdet
$u(3{,}25) \approx 4{,}06\,\text{V}$.
\sv{e} Kurvan startar i $3\,\text{V}$, har en dal vid $(0{,}75;\, 2{,}49)$ och en topp vid
$(3{,}25;\, 4{,}06)$ och faller sedan.

\svar{2.1} \sv{a} $i'(t) = 200\cos(50t)$
\sv{b} $50t = \pi/2 + k\pi$, alltså $t = \pi/100 + k\pi/50$,
$k = 0, 1, 2, \ldots$; de två första är $t_1 \approx 0{,}031\,\text{s}$ och
$t_2 \approx 0{,}094\,\text{s}$.
\sv{c} $i''(t) = -10\,000\sin(50t)$: $i''(t_1) = -10\,000 < 0$ ger maximum,
$i''(t_2) = 10\,000 > 0$ ger minimum.
\sv{d} $4\,\text{A}$ och $-4\,\text{A}$, vilket stämmer med amplituden $4\,\text{A}$.

\svar{2.2} \sv{a} $u'(t) = -4{,}8e^{-0{,}4t}$ \sv{b} $u'(2) \approx -2{,}16\,\text{V/s}$
\sv{c} $u''(2) \approx 0{,}86\,\text{V/s}^2 > 0$: förändringshastigheten ökar, det vill säga
spänningen sjunker allt långsammare.

\svar{2.3} \sv{a} $x = (10^{0{,}2} - 1)/2 \approx 0{,}29$
\sv{b} $G'(x) = 20/\bigl((2x + 1)\ln 10\bigr) \approx 5{,}48$ dB per enhet

\svar{3.1} \sv{a} $3\cos(3x)$ \sv{b} $-5\sin(5x)$ \sv{c} $4\cos(x)$ \sv{d} $4\sin(2x)$
\sv{e} $\cos(x) - \sin(x)$

\svar{3.2} \sv{a} $5e^{5x}$ \sv{b} $-6e^{-2x}$ \sv{c} $e^x + 2x$ \sv{d} $2^x \ln 2$
\sv{e} $-2{,}5e^{-x/4}$

\svar{3.3} \sv{a} $\frac{1}{x}$ \sv{b} $\frac{1}{x}$ \sv{c} $\frac{3}{x}$
\sv{d} $\frac{1}{x \ln 10}$ \sv{e} $\frac{1}{x} + e^x$

\svar{3.4} \sv{a} $6e^{3x} - 4\cos(4x)$ \sv{b} $2x + \frac{1}{x}$ \sv{c} $-5\sin(x) + 3e^{-x}$
\sv{d} $8\cos(2x) - 8\sin(2x)$

\svar{3.5} \sv{a} $2$ \sv{b} $0$ \sv{c} $0{,}5$ \sv{d} $0$ \sv{e} $-3e^{-1} \approx -1{,}10$

\svar{3.6} \sv{a} $u'(t) = 1000\pi\cos(100\pi t)$
\sv{b} $i_C(t) \approx 0{,}148\cos(100\pi t)\,\text{A}$ \sv{c} Cirka $148\,\text{mA}$
\sv{d} Vid $t = 0$, när spänningen passerar noll och ändras snabbast.

\svar{3.7} \sv{a} $i'(t) = 400\cos(200t)$ \sv{b} $u_L(t) = 10\cos(200t)\,\text{V}$
\sv{c} $10\,\text{V}$ \sv{d} Spänningen ligger $90^{\circ}$ före strömmen.

\svar{3.8} \sv{a} $u'(t) = -40e^{-2t}\,\text{V/s}$ \sv{b} $-40\,\text{V/s}$
\sv{c} $-40e^{-1} \approx -14{,}7\,\text{V/s}$
\sv{d} Derivatan är negativ, så spänningen minskar, snabbast i början.
\sv{e} $-u(t)/0{,}5 = -20e^{-2t}/0{,}5 = -40e^{-2t}$, samma uttryck som $u'(t)$ i a).

\svar{3.9} \sv{a} $f'(x) = e^{-2x}(1 - 2x)$ \sv{b} $x = 0{,}5$
\sv{c} Maximum: $f''(x) = e^{-2x}(4x - 4)$ ger $f''(0{,}5) = -2e^{-1} \approx -0{,}736 < 0$.
\sv{d} $f(0{,}5) = 0{,}5e^{-1} \approx 0{,}184$

\svar{3.10} \sv{a} $p'(t) = 100e^{-t}(1 - t)$ \sv{b} $t = 1\,\text{s}$
\sv{c} Maximum, eftersom $p''(1) = -100e^{-1} < 0$.
\sv{d} $p(1) = 100e^{-1} \approx 36{,}8\,\text{W}$

\svar{3.11} \sv{a} $G'(x) = \frac{20}{x \ln 10}$ \sv{b} $\approx 8{,}69$ dB per enhet
\sv{c} $\approx 0{,}869$ dB per enhet
\sv{d} Derivatan är tio gånger mindre vid $x = 10$: förstärkningen i dB växer allt långsammare,
eftersom den logaritmiska skalan komprimerar stora värden.

\svar{3.12} \sv{a} $1$ \sv{b} $-1$ \sv{c} $y = 1 - x$
\sv{d} Vid $x = 1$, det vill säga efter en tidskonstant ($\tau = 1$): hade urladdningen fortsatt
lika snabbt som i starten hade den varit klar efter tiden $\tau$.

\svar{3.13} \sv{a} Med $k = -1/\tau$ blir $u'(t) = -U_0e^{-t/\tau}/\tau = -u(t)/\tau$.
\sv{b} Förändringshastigheten är proportionell mot den aktuella spänningen: ju lägre spänning,
desto långsammare sjunker den.
\sv{c} $u'(0) = -5\,\text{V/s}$, $u'(\tau) \approx -1{,}84\,\text{V/s}$
\sv{d} $e^{-1} \approx 37\,\%$

\svar{3.14} \sv{a} Falskt: derivatan är $\cos x$. \sv{b} Falskt: den är $3e^{3x}$.
\sv{c} Falskt: $\ln(5x) = \ln 5 + \ln x$, så derivatan är $\frac{1}{x}$.
\sv{d} Sant: $\cos(kx)$ har derivatan $-k\sin(kx)$, här med $k = 2$.
\sv{e} Sant: $e^{kx}$ har derivatan $ke^{kx}$, och här är $k = 1$.
\sv{f} Sant: konstanten kan flyttas utanför, $\bigl(Cf(x)\bigr)' = Cf'(x)$.

\svar{3.15} \sv{a} Den inre faktorn $3$ ska med: $3\cos(3x)$.
\sv{b} Exponenten ändras inte, bara en faktor tillkommer: $-2e^{-2x}$.
\sv{c} Konstanten i logaritmen försvinner: $\frac{1}{x}$.
\sv{d} Derivatan av cosinus är minus sinus: $-4\sin(x)$.
\sv{e} Potensregeln gäller inte när $x$ står i exponenten: $2^x \ln 2$.

\svar[Testa dig själv]{T} \sv{1} $f'(x) = 2e^{2x} + \frac{1}{x}$
\sv{2} $f'(\pi/2) = \cos(\pi/2) = 0$

\end{facit}
